\begin{figure}[H]
\centering

\begin{center}
\textbf{Completeness Dependency Map}
\end{center}

\begin{tikzpicture}[>=Latex]

% Nodes
\node[depaxiom]   (comp)  at (0,0)       {Completeness Axiom};

\node[deptheorem] (nip)   at (-2.9,-1.7) {Nested Interval\\Property};
\node[deptheorem] (sqrt)  at ( 3.4,-1.7) {Existence of $\sqrt{a}$};

\node[deptheorem] (arch)  at (-2.9,-3.3) {Archimedean Property};
\node[deptheorem] (floor) at (-2.9,-4.9) {Floor Lemma};
\node[deptheorem] (densQ) at (-2.9,-6.5) {Density of $\mathbb{Q}$};
\node[deptheorem] (densI) at (-2.9,-8.1) {Density of the\\Irrationals in $\mathbb{R}$};

% Edges
\draw[depimplies] (comp) -- (nip);
\draw[depimplies] (comp) -- (sqrt);
\draw[depimplies] (nip)  -- (arch);
\draw[depimplies] (arch) -- (floor);
\draw[depimplies] (floor) -- (densQ);
\draw[depimplies] (densQ) -- (densI);

\end{tikzpicture}

\begin{center}
\DependencyFigureLegend
\end{center}

\caption{Dependency map for the completeness subchapter. The Completeness Axiom yields the Nested Interval Property and the existence of $\sqrt{a}$; the Nested Interval Property then supports the Archimedean Property, the Floor Lemma, and the density results for $\mathbb{Q}$ and for the irrational numbers.}
\label{fig:completeness-chain}
\end{figure}

\begin{remark*}[Connections forward]
Completeness appears in every major limit theorem in the notes.
The Monotone Convergence Theorem (convergence chapter) constructs the
limit of a bounded monotone sequence as $\sup\{a_n : n \in \mathbb{N}\}$;
existence of this supremum requires completeness.
The Bolzano--Weierstrass Theorem applies the Nested Interval Property
directly to extract a convergent subsequence from any bounded sequence.
The Cauchy Criterion identifies convergent sequences purely in terms of
their own spacing, without naming the limit in advance; its proof
constructs the limit as a supremum and again requires completeness.
In metric spaces (Distance chapter), completeness in the sense of
Cauchy sequences is the natural generalization of the same idea to
abstract settings.
Every $\varepsilon$-$\delta$ proof that locates a limit or extremum is
drawing on this chapter's infrastructure.
\end{remark*}